% !TeX root = ../Book1.tex
\begin{figure}[htbp]
	\centering
	\begin{tikzpicture}[scale=0.88, every node/.style={font=\small}]
		\begin{scope}
			\draw[thick] (-2,-1.4) .. controls (-2.5,-0.2) and (-2.1,1.6) .. (-0.8,1.8) .. controls (0.2,1.4) and (1.1,2.1) .. (1.9,1.3) .. controls (2.4,0.3) and (1.8,-1.6) .. (0.9,-1.8) .. controls (-0.3,-1.4) and (-1.2,-2) .. cycle;
			\draw[GeometryBlue,thick] (-1.2,-0.9) rectangle (1.2,0.9);
			\foreach \x in {-0.6,0,0.6}{\draw[GeometryBlue] (\x,-0.9)--(\x,0.9);}
			\foreach \y in {-0.45,0,0.45}{\draw[GeometryBlue] (-1.2,\y)--(1.2,\y);}
			\draw (-1.2,0.9)--(-0.8,1.8);
			\draw (1.2,-0.9)--(0.9,-1.8);
			\node at (0,2.25) {原区域 \(D\)};
			\node[fill=white,inner sep=2pt] at (0,-0.2) {\(K_n\)};
			\node at (-1.8,1.7) {\(C\)};
			\node at (0,-2.3) {紧集附近的网格逐次变细};
		\end{scope}
		\draw[->,thick] (2.7,0.2)--(3.7,0.2);
		\draw[<-,thick] (2.7,-0.2)--(3.7,-0.2);
		\begin{scope}[xshift=6.4cm]
			\draw[thick] (-1.8,-1.8) rectangle (1.8,1.8);
			\draw[GeometryBlue,thick] (-1.2,-0.9) rectangle (1.2,0.9);
			\foreach \x in {-0.6,0,0.6}{\draw[GeometryBlue] (\x,-0.9)--(\x,0.9);}
			\foreach \y in {-0.45,0,0.45}{\draw[GeometryBlue] (-1.2,\y)--(1.2,\y);}
			\draw (-1.2,0.9)--(-1.8,1.8);
			\draw (1.2,-0.9)--(1.8,-1.8);
			\node at (0,2.25) {目标正方形 \(Q\)};
			\node at (0,-2.3) {再加密像侧，并把新弦复制回去};
		\end{scope}
	\end{tikzpicture}
	\caption{双向网格控制的示意。每次只在内部添加有限条弦；原边界上的对应保持不变。两边都要加密，才能保证极限及其逆都连续。图中的光滑边界只为示意，证明允许任意 Jordan 曲线。}
	\label{g1:fig:double-mesh}
\end{figure}
